\documentclass[a4paper]{article}

\usepackage[margin=1in]{geometry} % reduce margin

% Basic packages
\usepackage[utf8]{inputenc}
\usepackage[T1]{fontenc}
\usepackage{textcomp}
\usepackage{url}
\usepackage{booktabs}
\usepackage{enumitem}
\usepackage[dvipsnames]{xcolor}
\usepackage{xifthen}

% Math packages
\usepackage{amsmath, amsfonts, mathtools, amsthm, amssymb}
\usepackage{mathrsfs}
\usepackage{cancel}
\usepackage{bm}
\usepackage{systeme}
\usepackage{stmaryrd} % for \lightning

% Math shortcuts
\newcommand\N{\ensuremath{\mathbb{N}}}
\newcommand\R{\ensuremath{\mathbb{R}}}
\newcommand\Z{\ensuremath{\mathbb{Z}}}
\renewcommand\O{\ensuremath{\emptyset}}
\newcommand\Q{\ensuremath{\mathbb{Q}}}
\newcommand\C{\ensuremath{\mathbb{C}}}
\newcommand\F{\ensuremath{\mathbb{F}}}
\DeclareMathOperator{\sgn}{sgn}
\DeclareMathOperator{\Ker}{ker}
\DeclareMathOperator{\im}{Im}

% Logic symbols
\let\svlim\lim\def\lim{\svlim\limits}
\let\implies\Rightarrow
\let\impliedby\Leftarrow
\let\iff\Leftrightarrow
\let\epsilon\varepsilon
\newcommand\contra{\scalebox{1.1}{$\lightning$}}

% Useful commands
\definecolor{correct}{HTML}{009900}
\newcommand\correct[2]{\ensuremath{\:}{\color{red}{#1}}\ensuremath{\to }{\color{correct}{#2}}\ensuremath{\:}}
\newcommand\green[1]{{\color{correct}{#1}}}

% Horizontal rule
\newcommand\hr{
    \noindent\rule[0.5ex]{\linewidth}{0.5pt}
}

% Simple theorem environments (without fancy boxes for homework)
\theoremstyle{definition}
\newtheorem{definition}{Definition}
\newtheorem{theorem}{Theorem}
\newtheorem{lemma}{Lemma}
\newtheorem{proposition}{Proposition}
\newtheorem{corollary}{Corollary}
\newtheorem*{remark}{Remark}
\newtheorem*{note}{Note}
\newtheorem*{example}{Example}

% Problem environment
\newcounter{problem}
\newenvironment{problem}[1][]
{
    \stepcounter{problem}
    \section*{Problem \theproblem\ifx\relax#1\relax\else: #1\fi}
}
{}

% Subproblem environment
\newcounter{subproblem}[problem]
\newenvironment{subproblem}[1][]
{
    \stepcounter{subproblem}
    \subsection*{(\alph{subproblem})\ifx\relax#1\relax\else\ #1\fi}
}
{}

% Solution environment
\newenvironment{solution}
{
    \noindent\textbf{Solution:}\\
}
{
    
}

% Headers
\usepackage{fancyhdr}
\pagestyle{fancy}
\fancyhf{}
\fancyhead[L]{Sherlock Zhang}
\fancyhead[C]{Math 4102 - Midterm Problem List}
\fancyhead[R]{ID: 518915}
\fancyfoot[C]{\thepage}

% Title info
\title{Math 4102 - Midterm Problem List}
\author{Sherlock Zhang}
\date{\today}

\begin{document}

\maketitle

% =============================================================================
% HOMEWORK PROBLEMS START HERE
% =============================================================================

\begin{problem}[Theorem 30.1]

State and prove Egoroff's theorem.

\end{problem}

\vspace{2em}
\begin{solution}

\textbf{30.1 \quad Theorem (Egoroff):} Let $A$ be a bounded measurable set in $\mathbb{R}$. Let $f_n$ and $f$ be functions defined on $A$ such that each $f_n$ is measurable and $f_n \to f$ almost everywhere on $A$. Then given any $\varepsilon > 0$, there exists a measurable subset $A_\varepsilon \subset A$ such that $m(A \setminus A_\varepsilon) < \varepsilon$ and $f_n \to f$ uniformly on $A_\varepsilon$.

\medskip

\textbf{Proof:} Given $\varepsilon > 0$ and positive integers $m$ and $n$, define the set
\[
E_{m,n} = \bigcap_{i=n}^{\infty} \{x \mid |f_i(x) - f(x)| < 1/m\}.
\]
Note that $E_{m,n}$ is measurable. If $U$ is the subset on which $f_n \to f$, then clearly for any $m$, 
\[
U \subset \bigcup_{n=1}^{\infty} E_{m,n} \subset A.
\]

Now $m(U) = m(A)$ so
\[
m\!\left(\bigcup_{n=1}^{\infty} E_{m,n}\right) = m(A).
\]
But $E_{m,n} \subset E_{m,n+1}$ for all $m$ and $n$, so
\[
\lim_{n \to \infty} m(A \setminus E_{m,n})
= \lim_{n \to \infty} [m(A) - m(E_{m,n})]
= m(A) - \lim_{n \to \infty} m(E_{m,n})
= m(A) - m\!\left(\bigcup_{n=1}^{\infty} E_{m,n}\right)
= 0.
\]

Thus for each $m$, there exists an integer $n_m$ such that
\[
m(A \setminus E_{m,n_m}) < \varepsilon / 2^m.
\]

Now let
\[
A_\varepsilon = \bigcap_{m=1}^{\infty} E_{m,n_m}.
\]
Then $A_\varepsilon$ is measurable and
\[
m(A \setminus A_\varepsilon)
= m\!\left(A \setminus \bigcap_{m=1}^{\infty} E_{m,n_m}\right)
= m\!\left(\bigcup_{m=1}^{\infty} (A \setminus E_{m,n_m})\right)
\le \sum_{m=1}^{\infty} m(A \setminus E_{m,n_m})
< \sum_{m=1}^{\infty} \varepsilon / 2^m
= \varepsilon.
\]

We now claim that $f_n \to f$ uniformly on $A_\varepsilon$ since given any $m$, there exists an $n_m$ such that for all $n > n_m$, $|f_n(x) - f(x)| < 1/m$ everywhere on $E_{m,n_m}$. But $A_\varepsilon \subset E_{m,n_m}$ for every $m$, so for all $n > n_m$, $|f_n(x) - f(x)| < 1/m$ everywhere on $A_\varepsilon$. But this is exactly uniform convergence. \qed

\end{solution}

\newpage
\begin{problem}[Theorem 30.3]

State and prove Lusin's theorem.

\end{problem}

\vspace{2em}
\begin{solution}
\textbf{30.3 \quad Theorem (Lusin):} If $f$ is a measurable function on a bounded measurable set $A$, then given any $\varepsilon > 0$, there exists a closed set $C_\varepsilon \subset A$ such that $m(A \setminus C_\varepsilon) < \varepsilon$ and $f$ is continuous on $C_\varepsilon$. (That is, the restriction of $f$ to $C_\varepsilon$ is continuous.)

\medskip

\textbf{Proof:} First, let 
$f = \sum_{k=1}^{n} a_k \chi_{E_k}$
be a simple function on $A$. Given $\varepsilon > 0$, for each $k$ there exists a closed set $C_k \subset E_k$ such that
\[
m(E_k \setminus C_k) < \varepsilon / n
\]
(see Corollary 13.3). Now $ C = \bigcup_{k=1}^{n} C_k $ is closed, $m(A \setminus C) < \varepsilon$ (Exercise 31.23) and $f$ is continuous on $C$ since $C$ is equal to the disjoint union of closed sets $C_k$ and $f$ is constant on each $C_k$. (Exercise 31.24).

\medskip

Now let $f$ be any measurable function. Since $f = f_+ - f_-$, we may assume $f$ is non-negative. Since $f$ is measurable, $ f = \lim_{n \to \infty} f_n $ for a sequence $\{f_n\}$ of simple functions by Theorem 19.4.

\medskip

By the first part of the proof, given $\varepsilon > 0$, there exists for each $n$ a closed set $C_n$ such that $m(A \setminus C_n) < \varepsilon / 2^{n+1}$ and $f_n$ is continuous on $C_n$.

\medskip

Let 
\[
C_0 = \bigcap_{n=1}^{\infty} C_n.
\]
Then $C_0$ is closed and
\[
m(A \setminus C_0)
= m\!\left(A \setminus \bigcap_{n=1}^{\infty} C_n\right)
= m\!\left(\bigcup_{n=1}^{\infty} (A \setminus C_n)\right)
\le \sum_{n=1}^{\infty} m(A \setminus C_n)
< \varepsilon \sum_{n=1}^{\infty} \frac{1}{2^{n+1}}
= \varepsilon/2.
\]

We would like to conclude that $f$ is continuous on $C_0$, but we cannot since we need something like uniform convergence for this. However, by Egoroff's Theorem, there is a measurable subset $C \subset C_0$ such that $m(C_0 \setminus C) < \varepsilon/4$ [and thus $m(A \setminus C) = m(A \setminus C_0) + m(C_0 \setminus C) < 3\varepsilon/4$], and $f_n \to f$ uniformly on $C$.

\medskip

Now each $f_n$ is continuous on $C \subset C_0$ so that $f$, being the uniform limit of the $f_n$, is continuous on $C$. If $C$ is closed we are through. If not, take closed set $F \subset C$ such that $m(C \setminus F) < \varepsilon/4$. \qed
\end{solution}

\newpage
\begin{problem}[Definition 32.3]
State the definition of a normed linear space and give both finite- and $ \infty $-dimensional examples.
\end{problem}

\vspace{2em}
\begin{solution}
\textbf{Definition.} A real \emph{normed linear space} is a vector space $V$ (with addition and scalar multiplication denoted by $\bar{u} + \bar{v}$, $a\bar{u}$, respectively) such that for each $\bar{u} \in V$, there is a unique real number $\|\bar{u}\|$ (called the norm or length of $\bar{u}$) satisfying:
\begin{enumerate}
    \item $\|\bar{u}\| \ge 0$ for all $\bar{u} \in V$;
    \item $\|\bar{u}\| = 0$ if and only if $\bar{u} = \bar{0}$;
    \item $\|a\bar{u}\| = |a|\,\|\bar{u}\|$ for all $a \in \mathbb{R}$, $\bar{u} \in V$;
    \item $\|\bar{u} + \bar{v}\| \le \|\bar{u}\| + \|\bar{v}\|$ for all $\bar{u}, \bar{v} \in V$.
\end{enumerate}

\medskip

\textbf{Examples.}

\textbf{(Finite-dimensional example)}  
For $(x_1,\dots,x_n) \in \mathbb{R}^n$,
\[
\|(x_1,\dots,x_n)\| = \sqrt{x_1^2 + x_2^2 + \cdots + x_n^2}.
\]

\textbf{($\infty$-dimensional example)}  
For $f \in C([0,1])$,
\[
\|f\| = \int_0^1 |f(x)|\,dx.
\]
\end{solution}

\newpage
\begin{problem}[Example 32.15]
State the definition of a real Hilbert space and give both finite- and $ \infty $-dimensional examples.
\end{problem}

\vspace{2em}
\begin{solution}
A real \emph{Hilbert space} is a real inner product space $V$ that is complete with respect to the norm $\|\bar{u}\| = \sqrt{\bar{u} \cdot \bar{u}}$ induced by the inner product.

\medskip

\textbf{Examples.}

\textbf{(Finite-dimensional example)}
$\mathbb{R}^n$ with the standard inner product
\[
\bar{u} \cdot \bar{v} = \sum_{i=1}^{n} u_i v_i
\]
is a Hilbert space. Completeness follows from the completeness of $\mathbb{R}$.

\textbf{($\infty$-dimensional example)}
$L^2[a,b]$ with the inner product
\[
f \cdot g = \int_{[a,b]} fg\,dm
\]
is a Hilbert space. Completeness is given by the Riesz--Fischer theorem (Theorem 33.11).
\end{solution}

\newpage
\begin{problem}
Determine whether $\{\cos kx:k=0,1,2,\ldots\}$ is an orthogonal set in
$L^2[0,\pi]$.
\end{problem}

\vspace{2em}
\begin{solution}
For $j,k=0,1,2,\ldots$, consider
\[
    \int_0^\pi \cos jx \cos kx\,dx.
\]
If $j\neq k$, then by the product-to-sum identity,
\[
    \cos jx\cos kx
    =
    \frac{1}{2}\bigl[\cos((j-k)x)+\cos((j+k)x)\bigr].
\]
Therefore,
\[
\begin{aligned}
    \int_0^\pi \cos jx\cos kx\,dx
    &=
    \frac{1}{2}\int_0^\pi \cos((j-k)x)\,dx
    +
    \frac{1}{2}\int_0^\pi \cos((j+k)x)\,dx \\[4pt]
    &=
    \frac{1}{2}
    \left[\frac{\sin((j-k)x)}{j-k}\right]_0^\pi
    +
    \frac{1}{2}
    \left[\frac{\sin((j+k)x)}{j+k}\right]_0^\pi \\[4pt]
    &=0.
\end{aligned}
\]
Also, none of the functions $\cos kx$ is the zero function in $L^2[0,\pi]$.
Hence
\[
    \{\cos kx:k=0,1,2,\ldots\}
\]
is an orthogonal set in $L^2[0,\pi]$.
\end{solution}

\newpage
\begin{problem}[Theorem 32.13]
    State and prove Cauchy's inequality
\end{problem}

\vspace{2em}
\begin{solution}
\textbf{32.13 Cauchy Inequality:} If $V$ is an inner product space with inner product $\bar{u} \cdot \bar{v}$, then for any $\bar{u}, \bar{v} \in V$,
\[
|\bar{u} \cdot \bar{v}| \le \sqrt{\bar{u} \cdot \bar{u}} \, \sqrt{\bar{v} \cdot \bar{v}}.
\]

\textbf{Proof:} If $\bar{u} = \bar{0}$ or $\bar{v} = \bar{0}$, then equality holds trivially. Otherwise, consider the vectors
\[
\bar{u}' = \frac{\bar{u}}{\sqrt{\bar{u} \cdot \bar{u}}}
\quad \text{and} \quad
\bar{v}' = \frac{\bar{v}}{\sqrt{\bar{v} \cdot \bar{v}}}.
\]

It is easy to verify that $\bar{u}' \cdot \bar{u}' = \bar{v}' \cdot \bar{v}' = 1$, so that by the lemma,
\[
|\bar{u}' \cdot \bar{v}'| \le 1.
\]
But
\[
|\bar{u}' \cdot \bar{v}'|
= \frac{|\bar{u} \cdot \bar{v}|}{\sqrt{\bar{u} \cdot \bar{u}} \, \sqrt{\bar{v} \cdot \bar{v}}}
\le 1.
\]
\qed
\end{solution}


\newpage
\begin{problem}[Theorem 33.9]
State and prove Minkowsky's inequality in $ \mathcal{L}^2 $.
\end{problem}

\vspace{2em}
\begin{solution}
\textbf{Theorem (Minkowski Inequality in $L^2$).}
For $f,g \in L^2(A)$,
\[
\|f+g\|_2 \le \|f\|_2 + \|g\|_2,
\]
i.e.,
\[
\left( \int_A (f+g)^2 \, dm \right)^{1/2}
\le
\left( \int_A f^2 \, dm \right)^{1/2}
+
\left( \int_A g^2 \, dm \right)^{1/2}.
\]

\textbf{Proof:}
We expand $\|f+g\|_2^2$:
\[
\|f+g\|_2^2
= \int_A (f+g)^2\,dm
= \int_A f^2\,dm + 2\int_A fg\,dm + \int_A g^2\,dm.
\]
By the Cauchy--Schwarz inequality (Theorem 33.5),
\[
\int_A fg\,dm \le \left|\int_A fg\,dm\right| \le \|f\|_2\,\|g\|_2.
\]
Therefore,
\[
\|f+g\|_2^2
\le \|f\|_2^2 + 2\|f\|_2\,\|g\|_2 + \|g\|_2^2
= \bigl(\|f\|_2 + \|g\|_2\bigr)^2.
\]
Taking square roots of both sides gives
\[
\|f+g\|_2 \le \|f\|_2 + \|g\|_2.
\]
\qed
\end{solution}

\newpage
\begin{problem}[Page 120]
Define a trigonometric polynomial and a trigonometric series
\end{problem}

\vspace{2em}
\begin{solution}
A \textbf{trigonometric polynomial} is a finite sum of the form
\[
    \sum_{k=0}^{n} \left(a_k \cos kx + b_k \sin kx\right),
\]
where $a_k,b_k \in \mathbb{R}$.

\vspace{2em}
\noindent A \textbf{trigonometric series} is an infinite series of the form
\[
    \sum_{k=0}^{\infty} \left(a_k \cos kx + b_k \sin kx\right),
\]
where $a_k,b_k \in \mathbb{R}$.
\end{solution}

\newpage
\begin{problem}[Theorem 36.3]
State and prove Bessel's inequality for an orthonormal set in $L^2[-\pi,\pi]$.
\end{problem}

\vspace{2em}
\begin{solution}
Let $\{u_k\}_{k=0}^{\infty}$ be an orthonormal set in $L^2[-\pi,\pi]$, and let
$f \in L^2[-\pi,\pi]$. Define
\[
    c_k = f \cdot u_k = \int_{[-\pi,\pi]} f u_k \, dm.
\]
Then Bessel's inequality states that
\[
    \sum_{k=0}^{\infty} c_k^2 \le \int_{[-\pi,\pi]} f^2 \, dm.
\]

\textbf{Proof.}
By Corollary 36.2, for every $n$,
\[
    \int_{[-\pi,\pi]} f^2 \, dm \ge \sum_{k=0}^{n} c_k^2.
\]
Since each $c_k^2 \ge 0$, the partial sums
\[
    \sum_{k=0}^{n} c_k^2
\]
are increasing and bounded above by $\int_{[-\pi,\pi]} f^2 \, dm$.
Therefore the infinite series converges and satisfies
\[
    \sum_{k=0}^{\infty} c_k^2 \le \int_{[-\pi,\pi]} f^2 \, dm.
\]
\end{solution}

\newpage
\begin{problem}[Theorem 36.4]
State and prove the Riemann-Lebesgue lemma for the Fourier series of $f \in L^2[-\pi,\pi]$.
\end{problem}

\vspace{2em}
\begin{solution}
Let $\{u_k\}_{k=0}^{\infty}$ be an orthonormal set in $L^2[-\pi,\pi]$, and let
$f \in L^2[-\pi,\pi]$. Define
\[
    c_k = f \cdot u_k = \int_{[-\pi,\pi]} f u_k \, dm.
\]
The Riemann-Lebesgue lemma states that
\[
    \lim_{k\to\infty} c_k = 0.
\]

\textbf{Proof.}
By Bessel's inequality,
\[
    \sum_{k=0}^{\infty} c_k^2 \le \int_{[-\pi,\pi]} f^2 \, dm.
\]
Thus $\sum_{k=0}^{\infty} c_k^2$ is a convergent series of nonnegative terms.
Therefore,
\[
    \lim_{k\to\infty} c_k^2 = 0.
\]
Hence,
\[
    \lim_{k\to\infty} c_k = 0.
\]
\end{solution}

\newpage
\begin{problem}[Lemma 37.4]
Give two formulas for the Dirichlet kernel $D_n(t)$.
\end{problem}

\vspace{2em}
\begin{solution}
The \textbf{Dirichlet kernel} $D_n(t)$ is
\[
D_n(t)
=
\frac12+\sum_{k=1}^{n}\cos(kt).
\]

\vspace{2em}
\noindent Equivalently, using Lemma 37.4,
\[
D_n(t)=
\begin{cases}
0, & t=2k\pi,\ k\in\mathbb{Z},\\[6pt]
\displaystyle
\frac{\sin\left(\left(n+\frac12\right)t\right)}
{2\sin\left(\frac{t}{2}\right)},
& \text{otherwise}.
\end{cases}
\]
\end{solution}

\newpage
\begin{problem}[Lemma 37.5]
Express the partial sum $s_n(x)$ of the Fourier series for a function
$f\in L^2[-\pi,\pi]$ using the Lebesgue integral and the Dirichlet kernel.
\end{problem}

\vspace{2em}
\begin{solution}
By Lemma 37.5, if $f\in L^2[-\pi,\pi]$, then the Fourier partial sum
$s_n(x)$ can be written as
\[
    s_n(x)
    =
    \frac{1}{\pi}\int_0^\pi
    \bigl[f(x+t)+f(x-t)\bigr]D_n(t)\,dt,
\]
where
\[
D_n(t)=
\begin{cases}
0, & t=2k\pi,\ k\in\mathbb{Z},\\[6pt]
\displaystyle
\frac{\sin\left(\left(n+\frac12\right)t\right)}
{2\sin\left(\frac{t}{2}\right)},
& \text{otherwise}.
\end{cases}
\]
Here $f$ is understood to be extended periodically with period $2\pi$.
\end{solution}

\newpage
\begin{problem}[Corollary 37.6]
Give a formula for the Fej\'er kernel $K_n(t)$ and relate it to the Dirichlet kernel.
\end{problem}

\vspace{2em}
\begin{solution}
The \textbf{Fej\'er kernel} $K_n(t)$ is
\[
K_n(t)=
\begin{cases}
0, & t=2m\pi,\ m\in\mathbb{Z},\\[6pt]
\displaystyle
\frac{\sin^2\left(\frac{nt}{2}\right)}
{2n\sin^2\left(\frac{t}{2}\right)},
& \text{otherwise}.
\end{cases}
\]

\vspace{2em}
\noindent Its relation with the Dirichlet kernel is
\[
    K_n(t)=\frac{1}{n}\sum_{k=0}^{n-1}D_k(t).
\]

Indeed, using the definition of $D_k(t)$,
\[
\frac{1}{n}\sum_{k=0}^{n-1}D_k(t)
=
\frac{1}{2n\sin(t/2)}
\sum_{k=0}^{n-1}
\sin\left(\left(k+\frac12\right)t\right).
\]
By Lemma 37.4,
\[
\sum_{k=0}^{n-1}
\sin\left(\left(k+\frac12\right)t\right)
=
\sum_{j=1}^{n}\sin\left((2j-1)\frac{t}{2}\right)
=
\frac{\sin^2(nt/2)}{\sin(t/2)}.
\]
Therefore,
\[
\frac{1}{n}\sum_{k=0}^{n-1}D_k(t)
=
\frac{\sin^2(nt/2)}
{2n\sin^2(t/2)}
=
K_n(t).
\]
\end{solution}

\newpage
\begin{problem}[Definition 37.1]
Define the Ces\`aro $(C,1)$-sum of a series $\sum_{k=1}^{\infty} a_k$. 
Give an example of a divergent series with a convergent $(C,1)$-sum.
\end{problem}

\vspace{2em}
\begin{solution}
A series $\sum_{k=1}^{\infty} a_k$ is said to have \textbf{$(C,1)$-sum} (or Ces\`aro sum) $s$ if
\[
\lim_{n\to\infty}\sigma_n = s,
\]
where
\[
\sigma_n = \frac{s_1+s_2+\cdots+s_n}{n}
\quad \text{and} \quad
s_n = \sum_{k=1}^{n} a_k.
\]
The sequence $\{\sigma_n\}$ is called the sequence of arithmetic means.

\vspace{2em}
\noindent An example is the series
\[
1 - 1 + 1 - 1 + \cdots,
\]
which diverges, but has $(C,1)$-sum $\frac{1}{2}$, since
\[
\sigma_n =
\begin{cases}
\frac{1}{2}, & \text{if } n \text{ is even},\\[6pt]
\frac{n+1}{2n}, & \text{if } n \text{ is odd},
\end{cases}
\]
and hence $\sigma_n \to \frac{1}{2}$ as $n\to\infty$.
\end{solution}

\newpage
\begin{problem}[Corollary 37.6]
Express the $(C,1)$-sum $\sigma_n(x)$ of the Fourier series for a function
$f\in L[-\pi,\pi]$ using the Lebesgue integral and the Fej\'er kernel.
\end{problem}

\vspace{2em}
\begin{solution}
Let
\[
    \sigma_n(x)=\frac{s_0(x)+s_1(x)+\cdots+s_{n-1}(x)}{n},
\]
where $s_n(x)$ is the $n$th Fourier partial sum. By Corollary 37.6,
\[
    \sigma_n(x)
    =
    \frac{1}{\pi}\int_0^\pi
    [f(x+t)+f(x-t)]K_n(t)\,dt,
\]
where $K_n(t)$ is the Fej\'er kernel of order $n$, given by
\[
K_n(t)=
\begin{cases}
0, & t=2m\pi,\ m\in\mathbb{Z},\\[6pt]
\displaystyle
\frac{\sin^2\left(\frac{nt}{2}\right)}
{2n\sin^2\left(\frac{t}{2}\right)},
& \text{otherwise}.
\end{cases}
\]
\end{solution}

\newpage
\begin{problem}[Theorem 37.8]
Prove that if $f\in L^2[-\pi,\pi]$ and $f$ is continuous, then
$\{\sigma_n\}$ converges uniformly to $f$.
\end{problem}

\vspace{2em}
\begin{solution}
Since $f$ is continuous on $[-\pi,\pi]$, it is bounded and uniformly continuous.
Thus there exists a number $M$ such that
\[
    |f(x)|<M
\]
for all $x\in[-\pi,\pi]$. By uniform continuity, given $\epsilon>0$, there exists
a $\delta>0$ such that $|x-y|<\delta$ in $[-\pi,\pi]$ implies
\[
    |f(x)-f(y)|<\frac{\epsilon}{2}.
\]
We will show that for some $N$, $n>N$ implies
\[
    |f(x)-\sigma_n(x)|<\epsilon
\]
for every $x\in[-\pi,\pi]$.

Now
\[
\begin{aligned}
|\sigma_n(x)-f(x)|
&\overset{37.6}{=}
\left|
\frac{2}{\pi}\int_0^\pi
\frac{f(x+t)+f(x-t)}{2}K_n(t)\,dt
-f(x)
\right| \\[6pt]
&\overset{37.7}{=}
\left|
\frac{2}{\pi}\int_0^\pi
\frac{f(x+t)+f(x-t)}{2}K_n(t)\,dt
-f(x)\frac{2}{\pi}\int_0^\pi K_n(t)\,dt
\right| \\[6pt]
&\le
\left|
\frac{2}{\pi}\int_0^\delta
\frac{f(x+t)+f(x-t)-2f(x)}{2}K_n(t)\,dt
\right| \\[6pt]
&\quad+
\left|
\frac{2}{\pi}\int_\delta^\pi
\frac{f(x+t)+f(x-t)-2f(x)}{2}K_n(t)\,dt
\right| \\[6pt]
&\le
\frac{\epsilon}{2}\frac{2}{\pi}\int_0^\delta K_n(t)\,dt \\[6pt]
&\quad+
\frac{2}{4n\pi\sin^2(\delta/2)}
\int_\delta^\pi
\left[|f(x+t)|+|f(x-t)|+2|f(x)|\right]dt \\[6pt]
&<
\frac{\epsilon}{2}
+
\frac{1}{2n\pi\sin^2(\delta/2)}
\int_\delta^\pi 4M\,dt.
\end{aligned}
\]

The second term is clearly less than $\epsilon/2$ for $n$ large enough, so the
result follows.
\end{solution}

\newpage
\begin{problem}[Theorem 40.2]
Prove that if $f\in L^2[-\pi,\pi]$ and $f$ is differentiable at $x$, then
$\{s_n(x)\}$ converges to $f(x)$.
\end{problem}

\vspace{2em}
\begin{solution}
Let $x_0\in[-\pi,\pi]$ and suppose $f$ is differentiable at $x_0$.
From Lemma 37.5, the $n$th Fourier partial sum is
\[
s_n(x_0)=\frac{1}{\pi}\int_0^\pi [f(x_0+t)+f(x_0-t)]D_n(t)\,dt.
\]

Also, for the particular case $f(x)\equiv 1$, we obtain
\[
1=\frac{2}{\pi}\int_0^\pi D_n(t)\,dt
\qquad \text{for } n=1,2,3,\dots
\]

Thus, for our general $f$, differentiable at $x_0$,
\[
s_n(x_0)-f(x_0)
=
\frac{1}{\pi}\int_0^\pi
[f(x_0+t)+f(x_0-t)-2f(x_0)]D_n(t)\,dt.
\]

Now define
\[
h(t)=\frac{f(x_0+t)+f(x_0-t)-2f(x_0)}{2\sin\frac{t}{2}}.
\]
Then $h(t)$ can be defined so as to be continuous at $t=0$. Therefore,
$h\in L^2[0,\pi]$. Thus,
\[
s_n(x_0)-f(x_0)
=
\frac{1}{\pi}\int_0^\pi
h(t)\sin\left(\left(n+\frac12\right)t\right)\,dt.
\]

Thus by the Riemann--Lebesgue Lemma (Corollary 36.4),
\[
s_n(x_0)-f(x_0)\to 0 \quad \text{as } n\to\infty.
\]
\end{solution}

\newpage
\begin{problem}[Theorem 41.2]
State and prove the Riemann localization lemma.
\end{problem}

\vspace{2em}
\begin{solution}
Let $f\in L[-\pi,\pi]$ and $x\in[-\pi,\pi]$, and let $\delta$ be any positive real number.
Then
\[
    \lim_{n\to\infty}
    \int_\delta^\pi [f(x+t)+f(x-t)]D_n(t)\,dt
    =
    0.
\]

\textbf{Proof.}
Consider
\[
g(t)=
\begin{cases}
0, & \text{on } [0,\delta),\\[6pt]
\displaystyle
\frac{f(x+t)+f(x-t)}{2\sin(t/2)},
& \text{on } [\delta,\pi].
\end{cases}
\]

Clearly $g\in L[-\pi,\pi]$ since
\[
    2\sin(t/2)\ge 2\sin(\delta/2)>0
    \qquad \text{on } [\delta,\pi].
\]
Thus
\[
\int_\delta^\pi [f(x+t)+f(x-t)]D_n(t)\,dt
=
\int_0^\pi g(t)\sin\left(n+\frac12\right)t\,dt.
\]

Therefore,
\[
\lim_{n\to\infty}
\int_\delta^\pi [f(x+t)+f(x-t)]D_n(t)\,dt
=
0
\]
by the Riemann--Lebesgue Lemma and Exercise 42.6.
\end{solution}

\newpage
\begin{problem}[Exercise 19]
Compute the complex exponential Fourier series of
\[
    f(x)=\sin x+2\cos 3x.
\]
\end{problem}

\vspace{2em}
\begin{solution}
We write the complex exponential Fourier series in the form
\[
    f(x)=\sum_{n=-\infty}^{\infty} c_n e^{inx}.
\]
Using Euler's formulas,
\[
    \sin x=\frac{e^{ix}-e^{-ix}}{2i},
    \qquad
    \cos 3x=\frac{e^{3ix}+e^{-3ix}}{2}.
\]
Therefore,
\[
\begin{aligned}
    f(x)
    &=\sin x+2\cos 3x \\[4pt]
    &=\frac{e^{ix}-e^{-ix}}{2i}
      +2\left(\frac{e^{3ix}+e^{-3ix}}{2}\right) \\[4pt]
    &=-\frac{i}{2}e^{ix}
      +\frac{i}{2}e^{-ix}
      +e^{3ix}
      +e^{-3ix}.
\end{aligned}
\]
Thus the only nonzero complex Fourier coefficients are
\[
    c_1=-\frac{i}{2},
    \qquad
    c_{-1}=\frac{i}{2},
    \qquad
    c_3=1,
    \qquad
    c_{-3}=1.
\]
All other coefficients are zero. Hence
\[
    f(x)
    =
    \frac{i}{2}e^{-ix}
    -\frac{i}{2}e^{ix}
    +e^{-3ix}
    +e^{3ix}.
\]
\end{solution}

\end{document}
